% !TEX root = ../proyect.tex

\chapter{Marco tecnológico}\label{cap:marco-tecnologico}

Este capítulo describe el conjunto de tecnologías, frameworks y paradigmas que conforman la base técnica del proyecto LinceUS Assistant. Se complementa con el Capítulo~\ref{cap:arquitectura}, donde se detallan la arquitectura del sistema, las decisiones de diseño y los diagramas de componentes y despliegue. Aquí se justifica la elección de cada componente del stack tecnológico.

\section{Tecnologías de procesamiento conversacional}\label{sec:procesamiento-conversacional}

\subsection{Rasa Framework}

Rasa es un framework de código abierto para la construcción de asistentes conversacionales basados en IA. A diferencia de soluciones comerciales como Dialogflow o Lex de AWS, Rasa ofrece:

\begin{itemize}
    \item \textbf{Control total sobre el modelo}: Los datos de entrenamiento y los modelos permanecen bajo control del desarrollador
    \item \textbf{Personalización completa}: Políticas de diálogo y pipelines de NLU totalmente configurables
    \item \textbf{Ejecución local}: No depende de servicios externos para funcionar
    \item \textbf{Multilenguaje}: Soporte nativo para español mediante modelos de spaCy
\end{itemize}

\subsubsection{Componentes de Rasa utilizados}

\textbf{Rasa NLU (Natural Language Understanding)}

Rasa NLU se encarga de interpretar la entrada del usuario y extraer:
\begin{itemize}
    \item \textbf{Intenciones (intents)}: ¿Qué quiere hacer el usuario? (e.g., \texttt{consultar\_horario}, \texttt{buscar\_asignatura})
    \item \textbf{Entidades (entities)}: Información específica mencionada por el usuario (e.g., nombre de asignatura, curso, centro)
\end{itemize}

\textbf{Rasa Core}

Rasa Core gestiona el flujo de conversación mediante:
\begin{itemize}
    \item \textbf{Stories}: Conversaciones de ejemplo que entrenan el modelo de diálogo
    \item \textbf{Rules}: Reglas determinísticas para situaciones específicas
    \item \textbf{Políticas}: Algoritmos de aprendizaje que deciden la siguiente acción del bot
\end{itemize}

\textbf{Rasa Actions}

Las acciones personalizadas permiten al bot ejecutar código Python para:
\begin{itemize}
    \item Consultar la base de datos
    \item Generar respuestas dinámicas
    \item Llamar a APIs externas
    \item Procesar lógica de negocio compleja
\end{itemize}

\subsection{Procesamiento de Lenguaje Natural}

\subsubsection{spaCy}

spaCy proporciona capacidades avanzadas de NLP en español:

\begin{itemize}
    \item \textbf{Tokenización}: Separación del texto en palabras y signos de puntuación
    \item \textbf{Lematización}: Reducción de palabras a su forma base
    \item \textbf{Análisis morfológico}: Identificación de categorías gramaticales
    \item \textbf{Vectores de palabras}: Representaciones numéricas del significado de las palabras
\end{itemize}

El modelo \texttt{es\_core\_news\_md} incluye vocabulario específico del español y ha sido entrenado en corpus de textos en español.

\subsubsection{Coincidencia difusa (Fuzzy Matching)}

La librería RapidFuzz implementa algoritmos de distancia de cadenas que permiten:

\begin{itemize}
    \item Tolerar errores ortográficos (e.g., "Fundamnetos de Programacion")
    \item Reconocer abreviaturas (e.g., "FP" por "Fundamentos de Programación")
    \item Manejar variaciones en la escritura (e.g., mayúsculas, tildes)
\end{itemize}

Se utiliza el algoritmo de \textit{Levenshtein distance} con un umbral de similitud del 80\% para considerar coincidencias válidas.

\section{Tecnologías de Inteligencia Artificial}\label{sec:ia}

\subsection{Ollama y modelos de lenguaje grandes (LLMs)}

Ollama es una plataforma que permite ejecutar modelos de lenguaje de gran tamaño localmente, sin necesidad de conexión a servicios cloud. En LinceUS Assistant se utiliza para dos propósitos principales:

\subsubsection{Text-to-SQL}

El sistema implementa un pipeline de generación de consultas SQL a partir de lenguaje natural:

\begin{enumerate}
    \item El usuario formula una pregunta en lenguaje natural
    \item Ollama (usando un modelo como Llama 3 o Mistral) genera la consulta SQL correspondiente
    \item La consulta se valida y ejecuta contra la base de datos
    \item Los resultados se procesan y formatean para presentarlos al usuario
\end{enumerate}

Este enfoque permite consultas complejas sin necesidad de definir explícitamente todas las posibles preguntas.

\subsubsection{Generación de respuestas contextualizadas}

Para consultas que no se ajustan a patrones predefinidos, el modelo LLM genera respuestas naturales basándose en:

\begin{itemize}
    \item El contexto de la conversación
    \item La información recuperada de la base de datos
    \item El conocimiento previo del modelo
\end{itemize}

\subsection{Embeddings y búsqueda semántica}

Para la recuperación de información de documentos (planes docentes, información de trámites), se utiliza:

\begin{itemize}
    \item \textbf{Modelos de embeddings}: Transforman texto en vectores numéricos que capturan el significado semántico
    \item \textbf{pgvector}: Extensión de PostgreSQL que permite almacenar y buscar vectores eficientemente
    \item \textbf{Búsqueda de similitud}: Encuentra documentos relevantes usando distancia coseno entre vectores
\end{itemize}

Este enfoque implementa un sistema RAG (\textit{Retrieval-Augmented Generation}) que mejora la precisión de las respuestas al combinar búsqueda semántica con generación de lenguaje.

\section{Tecnologías de base de datos}\label{sec:base-datos}

\subsection{PostgreSQL}

PostgreSQL es el sistema de gestión de bases de datos relacional utilizado. Se ha elegido por:

\begin{itemize}
    \item \textbf{Robustez}: Base de datos ACID con alta integridad de datos
    \item \textbf{Extensibilidad}: Soporte para tipos de datos personalizados y extensiones
    \item \textbf{Rendimiento}: Optimización de consultas complejas y soporte para índices avanzados
    \item \textbf{Código abierto}: Sin costes de licencia y comunidad activa
\end{itemize}

\subsection{Supabase}

Supabase proporciona una capa de servicios sobre PostgreSQL:

\begin{itemize}
    \item \textbf{API REST automática}: Generación automática de endpoints para operaciones CRUD
    \item \textbf{Real-time}: Suscripciones a cambios en la base de datos
    \item \textbf{Autenticación}: Sistema de usuarios y permisos integrado
    \item \textbf{Almacenamiento}: Para archivos y documentos
    \item \textbf{Funciones Edge}: Ejecución de código serverless
\end{itemize}

La API REST de Supabase simplifica la integración con el backend de Python, eliminando la necesidad de escribir queries SQL manualmente para operaciones básicas.

\subsection{pgvector}

La extensión pgvector añade capacidades de búsqueda vectorial a PostgreSQL:

\begin{itemize}
    \item Almacenamiento eficiente de vectores de alta dimensionalidad (hasta 16,000 dimensiones)
    \item Índices IVFFlat y HNSW para búsquedas rápidas de vecinos más cercanos
    \item Operadores de similitud (distancia euclidiana, producto punto, similitud coseno)
\end{itemize}

Esto permite implementar búsqueda semántica directamente en la base de datos, sin necesidad de sistemas externos como Pinecone o Weaviate.

\section{Tecnologías de frontend}\label{sec:frontend}

\subsection{Node.js}

El frontend del asistente está desarrollado usando Node.js, que proporciona:

\begin{itemize}
    \item \textbf{JavaScript del lado del servidor}: Mismo lenguaje en frontend y backend de la interfaz
    \item \textbf{NPM}: Gestor de paquetes con millones de librerías disponibles
    \item \textbf{Rendimiento}: Motor V8 de alta velocidad
\end{itemize}

\subsection{Framework web}

Se utiliza un framework minimalista que permite:

\begin{itemize}
    \item Servir la interfaz de chat web
    \item Gestionar la conexión WebSocket con el backend de Rasa
    \item Renderizar mensajes con formato enriquecido (tablas, enlaces, imágenes)
\end{itemize}

\section{Herramientas de desarrollo}\label{sec:herramientas-desarrollo-marco}

\subsection{Control de versiones}

\textbf{Git} es el sistema de control de versiones utilizado, con \textbf{GitHub} como plataforma de hosting. La estructura del repositorio incluye:

\begin{itemize}
    \item \texttt{/actions}: Acciones personalizadas de Rasa
    \item \texttt{/data}: Datos de entrenamiento (intents, stories, rules)
    \item \texttt{/models}: Modelos entrenados de Rasa
    \item \texttt{/frontend}: Código de la interfaz web
    \item \texttt{/docs}: Documentación del proyecto
\end{itemize}

\subsection{Gestión de dependencias}

Las dependencias de Python se gestionan mediante:

\begin{itemize}
    \item \textbf{requirements.txt}: Lista de todas las librerías necesarias con versiones específicas
    \item \textbf{Entorno virtual (venv)}: Aislamiento de dependencias del sistema
\end{itemize}

\subsection{Variables de entorno}

La librería \texttt{python-dotenv} gestiona la configuración mediante archivos \texttt{.env}:

\begin{itemize}
    \item Credenciales de Supabase (URL, API key)
    \item Configuración de Ollama (URL del servidor, modelo a utilizar)
    \item Parámetros de configuración del sistema
\end{itemize}

Esto permite cambiar la configuración sin modificar el código y evita exponer credenciales en el repositorio público.

\section{Consideraciones de seguridad}\label{sec:seguridad}

\subsection{Protección de datos}

\begin{itemize}
    \item \textbf{Credenciales}: Nunca se almacenan en el código, siempre en variables de entorno
    \item \textbf{API keys}: Se utilizan claves específicas con permisos mínimos necesarios
    \item \textbf{Conexiones}: HTTPS para todas las comunicaciones con servicios externos
\end{itemize}

\subsection{Validación de entrada}

\begin{itemize}
    \item Sanitización de consultas SQL para prevenir inyección SQL
    \item Validación de tipos de datos en entradas de usuario
    \item Limitación de longitud de mensajes
\end{itemize}

\subsection{Ejecución local de modelos}

El uso de Ollama para ejecutar modelos localmente aporta beneficios de privacidad:

\begin{itemize}
    \item Las consultas de los usuarios no se envían a servicios externos
    \item No hay dependencia de APIs comerciales que puedan cambiar términos de uso
    \item Control total sobre los datos procesados
\end{itemize}

\section{Stack tecnológico completo}\label{sec:stack-completo}

La siguiente tabla resume todas las tecnologías empleadas en el proyecto:

\begin{table}[h]
\centering
\begin{tabular}{|l|l|l|}
\hline
\textbf{Capa} & \textbf{Tecnología} & \textbf{Versión} \\ \hline
Frontend & Node.js & 18.x \\ \hline
Frontend & HTML/CSS/JavaScript & ES6+ \\ \hline
Backend conversacional & Python & 3.10 \\ \hline
Backend conversacional & Rasa & 3.6.21 \\ \hline
Backend conversacional & Rasa SDK & 3.6.2 \\ \hline
Backend conversacional & spaCy & 3.8.0 \\ \hline
IA y LLMs & Ollama & latest \\ \hline
IA y LLMs & Modelos LLM & Llama 3 / Mistral \\ \hline
Base de datos & PostgreSQL & 15.x \\ \hline
Base de datos & pgvector & 0.5.0 \\ \hline
Plataforma BD & Supabase & cloud \\ \hline
Librerías NLP & RapidFuzz & 3.0.0 \\ \hline
Librerías NLP & python-dotenv & 1.0.0 \\ \hline
Control versiones & Git \& GitHub & 2.x \\ \hline
\end{tabular}
\caption{Stack tecnológico completo del proyecto LinceUS Assistant.}
\label{tab:stack-tecnologico}
\end{table}
